\begin{proof}[Proof of \cref{obj:lem:unrestricted-cell-identification}]
Fix \(P\in\mathcal M^{+}_{n,d,q}\). We use the conditions in \cref{obj:def:unrestricted-arrival-model-class}.

\begin{enumerate}
\item By \cref{obj:synth_4}, the potential outcomes are binary. Linearity of integration for these bounded variables therefore gives
\[
\tau(P)
=E_P[Y(1)]-E_P[Y(0)]
=P(Y(1)=1)-P(Y(0)=1),
\]
using \cref{obj:def:ate-functional}. % lean: CausalSmith.Stat.MarRareqLogfrontier.ate_eq_potentialOutcomeMass

\item We establish, for every \((a,x,s)\in\mathcal J_d\), the identity
\[
c_{axs}(P)=P(A=a,X=x,S=s,Y=1).
\]
First suppose \(p_{axs}(P)>0\). By \cref{obj:synth_16,obj:ass:occupied-cell-arrival} and \(q>0\),
\[
P(A=a,X=x,S=s,R=1)
=p_{axs}(P)\rho_{axs}(P)
\ge q\,p_{axs}(P)>0.
\]
Conditional independence in \cref{obj:ass:arrival-mar}, applied within this occupied cell, yields the cross-multiplied identity
\[
\begin{aligned}
&P(A=a,X=x,S=s,R=1,Y=1)\,p_{axs}(P)\\
&\qquad =
P(A=a,X=x,S=s,R=1)\,
P(A=a,X=x,S=s,Y=1).
\end{aligned}
\]
Since \(Y\) is binary and the arrived-cell probability is positive, \cref{obj:def:cell-functional} gives
\[
\begin{aligned}
c_{axs}(P)
&=p_{axs}(P)
\frac{P(A=a,X=x,S=s,R=1,Y=1)}
     {P(A=a,X=x,S=s,R=1)}\\
&=\frac{P(A=a,X=x,S=s,R=1,Y=1)\,p_{axs}(P)}
        {P(A=a,X=x,S=s,R=1)}\\
&=P(A=a,X=x,S=s,Y=1).
\end{aligned}
\]

In the remaining case, nonnegativity implies \(p_{axs}(P)=0\). Monotonicity of probability gives
\[
\begin{aligned}
0&\le P(A=a,X=x,S=s,Y=1)\le p_{axs}(P)=0,\\
0&\le P(A=a,X=x,S=s,R=1)\le p_{axs}(P)=0.
\end{aligned}
\]
Thus both probabilities vanish, and the zero-arrival convention in
\cref{obj:def:cell-functional} gives \(c_{axs}(P)=0\), proving the identity also in this case. % lean: CausalSmith.Stat.MarRareqLogfrontier.cellContribution_eq_cellOutcomeMass

\item For each \(a\in\{0,1\}\), define the finite measure obtained by restricting \(P\) to the arm-\(a\), outcome-one event:
\[
\nu_a^{\mathrm{out}}(\,\cdot\,)
:=P\bigl(\,\cdot\,\cap\{A=a,Y=1\}\bigr).
\]
The events \(\{X=x,S=s\}\), with \(x\in\{0,\ldots,d-1\}\) and \(s\in\{0,1\}\), are disjoint and cover the full-data record space. Finite additivity of this restricted measure therefore yields
\[
\begin{aligned}
&\sum_{x=0}^{d-1}\sum_{s\in\{0,1\}}
P(A=a,X=x,S=s,Y=1)\\
&\qquad =
\sum_{x=0}^{d-1}\sum_{s\in\{0,1\}}
\nu_a^{\mathrm{out}}(X=x,S=s)
=P(A=a,Y=1).
\end{aligned}
\] % lean: CausalSmith.Stat.MarRareqLogfrontier.sum_cellOutcomeMass_eq_armOutcomeMass

\item Expanding \cref{obj:def:cell-functional} over the two treatment arms, using \(g(1)=1\) and \(g(0)=-1\) from \cref{obj:synth_19}, and substituting the preceding identities gives
\[
\begin{aligned}
\Phi(P)
&=2\sum_{x=0}^{d-1}\sum_{s\in\{0,1\}}
\bigl(c_{1xs}(P)-c_{0xs}(P)\bigr)\\
&=2\bigl(P(A=1,Y=1)-P(A=0,Y=1)\bigr).
\end{aligned}
\] % lean: CausalSmith.Stat.MarRareqLogfrontier.cellFunctional_eq_armOutcomeMass

\item Fix \(a\in\{0,1\}\). Projecting the independence asserted in
\cref{obj:ass:randomized-independence} onto the selected potential outcome gives
\[
A\perp\!\!\!\perp Y(a),
\qquad
P(A=a,Y(a)=1)=P(A=a)\,P(Y(a)=1).
\]
By \cref{obj:ass:balanced-randomization},
\[
P(A=1)=\frac12,
\qquad
P(A=0)=1-P(A=1)=\frac12.
\]
The assigned-arm outcome in \cref{obj:synth_4} satisfies \(Y=Y(a)\) on \(\{A=a\}\). Consequently,
\[
\begin{aligned}
2P(A=a,Y=1)
&=2P(A=a,Y(a)=1)\\
&=2P(A=a)\,P(Y(a)=1)
=P(Y(a)=1).
\end{aligned}
\]
Applying this identity first to \(a=1\) and then to \(a=0\) gives
\[
2P(A=1,Y=1)=P(Y(1)=1),
\qquad
2P(A=0,Y=1)=P(Y(0)=1).
\] % lean: CausalSmith.Stat.MarRareqLogfrontier.randomized_armOutcomeMass

\item Subtracting the arm-zero identity from the arm-one identity and substituting into the expressions for \(\tau(P)\) and \(\Phi(P)\) proves
\[
\tau(P)
=\Phi(P)
=2\sum_{(a,x,s)\in\mathcal J_d}g(a)c_{axs}(P),
\]
where the final equality is \cref{obj:def:cell-functional}. % lean: CausalSmith.Stat.MarRareqLogfrontier.unrestricted_cell_identification
\end{enumerate}
\end{proof}
